\chapter{Barisan dan Limit}

Sebuah program kecil mencetak satu bilangan per baris. Aturannya hanya satu: suku
berikutnya adalah akar dari dua ditambah suku sebelumnya, dimulai dari satu. Yang
muncul di layar: 1.000000, lalu 1.732051, 1.931852, 1.982890, 1.995718, 1.998929,
dan seterusnya. Pada baris kedua belas layar menunjukkan 1.999999738545. Angka itu
tampak berhenti bergerak. Tetapi tampak dan benar adalah dua hal berbeda: program
tidak tahu apakah baris keseribu masih akan bergeser, dan ia juga tidak tahu ke
bilangan mana deretan ini menuju.

Contoh kedua lebih nakal. Barisan $(1+1/n)^n$ untuk $n=10$ bernilai 2.5937424601,
untuk $n=100$ bernilai 2.7048138294, untuk $n=1000$ bernilai 2.7169239322, dan
untuk $n=10^7$ bernilai 2.7182816941. Enam angka pertama tampak mantap pada 2.71828.
Tetapi nilai sebenarnya pada $n=10^7$ adalah 2.7182816925, sehingga tujuh digit
terakhir pada keluaran mesin itu adalah galat pembulatan, bukan tanda kedekatan.
Tabel dapat menipu ke dua arah sekaligus: membuat barisan yang divergen tampak
tenang, dan membuat barisan yang konvergen tampak melenceng. Bab 22 nanti
menjelaskan mengapa aritmetika mesin berperilaku begitu; bab ini menyiapkan alat
yang membuat kita tidak perlu bergantung pada tabel.

Bab 3 sudah memberi bahan bakunya: aksioma lapangan, aksioma urutan, dan sifat
supremum. Tetapi supremum hanya bekerja setelah himpunan sukunya terkumpul; ia
tidak memberi tahu apakah barisan berhenti bergerak. Yang kita perlukan adalah
definisi yang bisa diperiksa dengan kertas, lalu teorema yang mengubah pemeriksaan
itu menjadi pekerjaan terkira: definisi $\varepsilon$-$N$, sifat aljabar limit,
teorema apit, barisan monoton, subbarisan, kriteria Cauchy, dan limit superior.

\begin{untukanak}
Pernahkah kalian berjalan menuju pintu dengan aturan aneh: setiap langkah hanya
menempuh separuh dari jarak yang tersisa? Langkah pertama membawa kalian separuh
jalan. Langkah kedua memotong separuh dari sisa itu, jadi kalian sudah di tiga
perempat jalan. Langkah ketiga menambah seperdelapan lagi. Setiap langkah pasti
maju, tetapi selalu ada sisa yang belum terjangkau. Kalau aturan itu dijalankan
terus, kalian tidak pernah benar-benar menyentuh pintu.

Meski begitu, ada pertanyaan yang bisa dijawab dengan pasti. Berapa langkah yang
diperlukan supaya jarak ke pintu kurang dari sejengkal? Jawabannya bisa dihitung
sebelum melangkah, dan tidak berubah karena ruangan menjadi dua kali lebih panjang:
kalian hanya perlu beberapa langkah tambahan. Jadi ada dua hal berbeda di sini.
"Tidak pernah sampai" dan "sedekat apa pun bisa didekati" bukan dua pernyataan
yang bertentangan. Justru keduanya benar sekaligus.

Sekarang pertanyaan yang membuat penasaran. Bagaimana kalau aturannya diubah
sedikit, misalnya setiap langkah menempuh sepertiga sisa jarak, atau separuh
ditambah sedikit? Apakah kalian masih bisa menebak di mana perjalanan itu berakhir,
dan apakah kecepatannya masih bisa dihitung? Di bab ini kita akan tahu mengapa
pertanyaan semacam itu punya jawaban yang pasti, dan mengapa jawabannya tidak
pernah kita peroleh hanya dengan menonton angka di layar.
\end{untukanak}

\section{Definisi limit dan cara memakainya}
\tingkat{dasar}

\kamus{barisan}{sequence}{deretan bilangan yang diberi nomor urut, seperti nomor
rumah di satu sisi jalan yang tidak pernah habis}

Limit menjawab satu pertanyaan: seberapa dekat suku-suku akhir berkumpul, dan mulai
kapan. Dalam bahasa sehari-hari, barisan menuju $L$ bila setelah beberapa suku awal
dibuang, semua sisa sukunya sedekat apa pun yang diminta kepada $L$.

\begin{definisi}[Kekonvergenan]
Barisan $(x_n)$ bilangan real \emph{konvergen} ke $L \in \mathbb{R}$ jika untuk
setiap $\varepsilon>0$ terdapat $N \in \mathbb{N}$ sehingga untuk setiap $n>N$
berlaku $|x_n-L|<\varepsilon$. Kita tulis $x_n \to L$, dan $L$ disebut
\emph{limit}\index{limit barisan} barisan itu. Barisan yang tidak konvergen ke
bilangan real mana pun disebut \emph{divergen}\index{barisan divergen}.
\end{definisi}

Urutan kuantor pada definisi itu menentukan segalanya. Nilai $\varepsilon$ dipilih
lebih dahulu, lalu $N$ menyusul dan boleh bergantung pada $\varepsilon$. Membalik
urutan itu menghasilkan pernyataan yang jauh lebih kuat dan biasanya salah: ada
barisan yang memenuhi "untuk setiap $n$ ada $N$" secara hampa, karena $N=n$ selalu
bisa dipakai. Cara membaca yang sehat adalah sebagai permainan dua pihak. Penantang
menyebut satu $\varepsilon$, sekecil apa pun. Penjawab harus membalas dengan satu
$N$ yang berlaku untuk semua suku setelahnya. Barisan konvergen bila penjawab
selalu punya balasan, untuk setiap tantangan.

Gambar~\ref{fig:pita-epsilon} memotret permainan itu untuk $x_n=1+1/n$. Dua pita
berpusat di $L=1$ digambar dengan lebar berbeda, dan masing-masing punya ambangnya
sendiri: $N=2$ untuk $\varepsilon=0.5$, $N=10$ untuk $\varepsilon=0.1$. Tidak ada
yang berubah pada barisannya; yang berubah hanya permintaan penantang.

\begin{figure}[htbp]
\centering
\begin{tikzpicture}[x=0.72cm,y=3.0cm,>=Stealth,font=\footnotesize]
  \fill[layerteal!14] (0,0.5) rectangle (14,1.5);
  \fill[layerorange!20] (0,0.9) rectangle (14,1.1);
  \draw[->,ink] (-0.1,0.35) -- (14.6,0.35) node[right] {$n$};
  \draw[->,ink] (0,0.35) -- (0,2.15) node[above] {$x_n$};
  \draw[dashed,layergrey,thick] (0,1) -- (14.1,1);
  \node[layergrey,anchor=south west] at (0.08,1.02) {$L=1$};
  \foreach \n/\y in {1/2.0000,2/1.5000,3/1.3333,4/1.2500,5/1.2000,6/1.1667,%
      7/1.1429,8/1.1250,9/1.1111,10/1.1000,11/1.0909,12/1.0833,13/1.0769,14/1.0714}
    \fill[accent] (\n,\y) circle (2.2pt);
  \draw[dashed,layerteal,thick] (2,0.35) -- (2,1.52);
  \draw[dashed,layerorange,thick] (10,0.35) -- (10,1.12);
  \node[anchor=north] at (2,0.30) {$N_1=2$};
  \node[anchor=north] at (10,0.30) {$N_2=10$};
  \node[anchor=north west,align=left] at (4.6,2.28)
    {pita luar: $\varepsilon=0.5$\\ pita dalam: $\varepsilon=0.1$};
\end{tikzpicture}
\caption{Barisan $x_n=1+1/n$ masuk ke pita $\varepsilon$ di sekitar $L=1$. Dengan
$\varepsilon=0.5$ semua suku sesudah $n=2$ sudah di dalam pita; dengan
$\varepsilon=0.1$ ambangnya bergeser ke $n=10$. Pita yang lebih sempit menuntut
ambang yang lebih jauh.}
\label{fig:pita-epsilon}
\end{figure}

\begin{bacakeras}
Dua suara bergantian. Suara pertama bertanya, suara kedua menjawab. Bacakan pelan,
beri jeda pada setiap titik, dan jangan menambah kata di antaranya.

``Semakin lama?'' \ ``Semakin dekat.'' \ ``Seberapa dekat?'' \ ``Sedekat yang
diminta.'' \ ``Kalau yang diminta seperseribu?'' \ ``Tunggu seribu langkah, lalu
ukur.'' \ ``Kalau seperseratus juta?'' \ ``Tunggu seratus juta langkah. Sabar tidak
apa-apa.'' \ ``Jadi akan sampai?'' \ ``Tidak. Tidak pernah sampai.'' \ ``Lalu apa
gunanya menunggu?'' \ ``Sebab jaraknya mengecil, dan yang meminta bebas memilih
seberapa kecil.'' \ ``Dan kalau permintaannya berhenti?'' \ ``Maka tidak ada yang
perlu dibuktikan.''

Baris yang diulang tiga kali, makin pelan setiap kali: dekat bukan sampai; dekat
adalah jarak yang boleh kita tentukan sendiri.
\end{bacakeras}

Untuk $x_n=1/n$, balasannya langsung: ambil $N$ bilangan asli dengan $N>1/\varepsilon$,
yang ada menurut sifat Archimedes di Bab 3. Untuk $n>N$ berlaku
$|1/n-0|=1/n<1/N<\varepsilon$. Bila penantang meminta $\varepsilon=10^{-3}$,
jawabannya $N=1000$; pada $n=1000$ galatnya tepat $10^{-3}$ dan belum memenuhi
syarat ketat, sedangkan pada $n=1001$ galatnya $0.000999$. Perbedaan satu suku ini
bukan keisengan: definisi memakai $n>N$, bukan $n \ge N$, sehingga $N$ yang dipilih
harus benar-benar melewati ambang.

\begin{tangan}
Bahan: lakban, meteran, kapur, dan satu lembar catatan. Bentuk garis bilangan di
lantai sepanjang dua meter: nol di ujung kiri, dua di ujung kanan, dan angka satu
tepat di tengah. Tempel dua pita mendatar selebar telapak kaki, berpusat di angka
satu: pita lebar dari setengah sampai satu setengah meter, pita sempit dari sembilan
persepuluh sampai satu koma satu meter.

Sekarang berjalanlah ke posisi satu ditambah satu per n, untuk n sama dengan satu,
dua, tiga, sampai empat belas. Setiap kali, catat apakah tumit masuk pita lebar,
pita sempit, atau di luar keduanya, lalu berjalan mundur supaya tanda tidak
terinjak.

Yang harus terlihat: dua langkah pertama masih di luar pita sempit, langkah
kesepuluh berhenti tepat di tepinya, dan sesudah itu tidak ada lagi tumit yang
keluar. Pita yang lebih sempit menuntut langkah yang lebih banyak, sedangkan angka
satu tidak pernah tersentuh tumit.
\end{tangan}

Kegunaan definisi justru paling terasa ketika kita ingin membuktikan bahwa sesuatu
\emph{tidak} konvergen. Negasi definisi berbunyi: ada $\varepsilon>0$ sehingga untuk
setiap $N$ terdapat $n>N$ dengan $|x_n-L| \ge \varepsilon$. Perhatikan bahwa $L$
masih nama bebas di situ; untuk menyangkal kekonvergenan kita harus melakukannya
untuk setiap $L$.

\begin{contoh}
Barisan $x_n=(-1)^n$ divergen.
\end{contoh}

\begin{proof}[Bukti]
Andaikan $x_n \to L$ untuk suatu $L$. Ambil $\varepsilon=1$; ada $N$ sehingga
$|x_n-L|<1$ untuk semua $n>N$. Pilih $n>N$ yang genap; maka $x_n=1$ dan
$|1-L|<1$. Suku berikutnya $n+1$ juga melebihi $N$ dan ganjil, jadi $x_{n+1}=-1$
dan $|-1-L|<1$. Ketaksamaan segitiga memberi
\[
2 = |1-(-1)| \le |1-L| + |L-(-1)| < 1+1 = 2,
\]
suatu kemustahilan. Jadi tidak ada $L$ yang memenuhi definisi.
\end{proof}

Bukti tadi memperlihatkan pola yang akan sering muncul: untuk menyangkal, kita
memilih $\varepsilon$ yang lebih kecil daripada jarak antar dua kelompok suku.
Kelompok genap dan ganjil berjarak tetap $2$, sehingga tidak ada satu $N$ pun yang
dapat memisahkan keduanya bila $\varepsilon<1$.

Sifat pertama yang patut dicatat adalah ketunggalan. Kalimat sehari-harinya: sebuah
barisan tidak dapat berhenti di dua tempat berbeda sekaligus.

\begin{teorema}[Ketunggalan limit]
Jika $x_n \to L$ dan $x_n \to M$, maka $L=M$.
\end{teorema}

\begin{proof}[Bukti]
Andaikan $L \neq M$ dan misalkan $\varepsilon=|L-M|/2>0$. Ada $N_1$ dengan
$|x_n-L|<\varepsilon$ untuk $n>N_1$, dan $N_2$ dengan $|x_n-M|<\varepsilon$ untuk
$n>N_2$. Untuk $n>\max(N_1,N_2)$ kedua ketaksamaan berlaku sekaligus, sehingga
\[
|L-M| \le |L-x_n| + |x_n-M| < 2\varepsilon = |L-M| .
\]
Ketaksamaan $|L-M|<|L-M|$ mustahil, jadi pengandaian salah.
\end{proof}

Frasa "untuk $n$ yang cukup besar" berarti: ada $N$ sehingga sifat yang disebut
berlaku untuk semua $n>N$. Dengan kosakata itu, definisi limit ringkasnya: $x_n$
sedekat apa pun ke $L$ bila $n$ cukup besar.

Sifat kedua: barisan yang konvergen tidak dapat lari ke tak hingga.

\begin{teorema}
Setiap barisan konvergen terbatas.
\end{teorema}

\begin{proof}[Bukti]
Misalkan $x_n \to L$. Ambil $\varepsilon=1$; ada $N$ dengan $|x_n-L|<1$ untuk
$n>N$, sehingga $|x_n| \le |x_n-L|+|L| < |L|+1$. Hanya suku pertama sampai suku
ke-$N$ yang belum terkendali, dan jumlahnya berhingga. Maka
\[
M=\max\{|x_1|,|x_2|,\dots,|x_N|,|L|+1\}
\]
adalah batas yang berlaku untuk semua $n$.
\end{proof}

Kebalikannya tidak berlaku: $(-1)^n$ terbatas tetapi divergen. Keterbatasan adalah
syarat perlu, bukan syarat cukup; dari tabel yang tampak tenang kita hanya boleh
menyimpulkan keterbatasan.

\section{Sifat aljabar limit}
\tingkat{menengah}

Kalimat sehari-harinya: limit sopan terhadap penjumlahan, perkalian, dan pembagian,
selama pembaginya tidak menuju nol.

Sebelum membuktikan sifat-sifat itu, kita perlu satu kendali kecil: barisan yang
menuju bilangan tak nol akhirnya menjauh dari nol.

\begin{lemma}
Jika $y_n \to y$ dengan $y \neq 0$, maka ada $N$ sehingga $|y_n|>|y|/2$ untuk
semua $n>N$.
\end{lemma}

\begin{proof}[Bukti]
Ambil $\varepsilon=|y|/2>0$. Ada $N$ dengan $|y_n-y|<|y|/2$ untuk $n>N$. Maka
$|y_n| \ge |y|-|y_n-y| > |y|-|y|/2 = |y|/2$.
\end{proof}

Konsekuensinya untuk berhitung: limit dapat dihitung suku demi suku, dan satu-satunya
pengecualian adalah pembagian dengan nol.

\begin{teorema}[Sifat aljabar limit]
Jika $x_n \to x$ dan $y_n \to y$, maka
$x_n+y_n \to x+y$, \ $x_n y_n \to xy$, dan $c x_n \to cx$ untuk setiap
$c \in \mathbb{R}$. Jika $y \neq 0$ dan $y_n \neq 0$ untuk semua $n$, maka
$x_n/y_n \to x/y$.
\end{teorema}

\begin{proof}[Bukti]
Untuk jumlah, ambil $\varepsilon>0$. Pilih $N_1$ dengan $|x_n-x|<\varepsilon/2$
untuk $n>N_1$ dan $N_2$ dengan $|y_n-y|<\varepsilon/2$ untuk $n>N_2$. Untuk
$n>\max(N_1,N_2)$,
$|(x_n+y_n)-(x+y)| \le |x_n-x|+|y_n-y| < \varepsilon$. Kasus kelipatan skalar
serupa dan lebih singkat: untuk $c=0$ barisan konstan; untuk $c \neq 0$ pakai
$\varepsilon/|c|$.

Untuk perkalian, tulis selisihnya sebagai dua bagian,
\[
x_n y_n - xy = x_n(y_n-y) + (x_n-x)y .
\]
Karena $(x_n)$ konvergen, ia terbatas: ada $M>0$ dengan $|x_n| \le M$ untuk semua
$n$. Pilih $N$ sehingga untuk $n>N$ berlaku $|y_n-y|<\varepsilon/(2M)$ dan
$|x_n-x|<\varepsilon/(2(|y|+1))$. Maka
\[
|x_n y_n - xy| \le M|y_n-y| + |x_n-x|\,|y| < \frac{\varepsilon}{2}
+ \frac{\varepsilon}{2} = \varepsilon .
\]
Penyebut $|y|+1$ dipakai supaya tidak pernah nol ketika $y=0$; ia tetap konstanta
positif, jadi tidak mengubah kesimpulan.

Untuk hasil bagi, cukup ditunjukkan $1/y_n \to 1/y$, lalu gunakan hasil perkalian.
Menurut lemma, untuk $n>N_1$ berlaku $|y_n|>|y|/2$, sehingga
\[
\left|\frac{1}{y_n}-\frac{1}{y}\right| = \frac{|y-y_n|}{|y_n y|}
\le \frac{2}{|y|^2}\,|y_n-y| .
\]
Ruas kanan lebih kecil dari $\varepsilon$ asalkan $|y_n-y|<\varepsilon|y|^2/2$,
dan itu dapat dijamin dengan memilih $N_2$.
\end{proof}

Perhatikan teknik yang baru muncul: menambah dan mengurangkan suku yang sama untuk
memecah selisih menjadi dua bagian kecil. Hampir setiap bukti kekonvergenan yang
melibatkan perkalian memakainya, dan ambang yang diambil selalu yang terbesar.

Sifat berikutnya menyangkut urutan, bukan operasi. Kalimat sehari-harinya: limit
tidak dapat melompati pagar yang dipatuhi semua suku.

\begin{teorema}[Sifat urutan]
Jika $x_n \to x$ dan $x_n \ge 0$ untuk semua $n$, maka $x \ge 0$. Lebih umum,
jika $x_n \le y_n$ untuk semua $n$ yang cukup besar, $x_n \to x$, dan $y_n \to y$,
maka $x \le y$.
\end{teorema}

\begin{proof}[Bukti]
Andaikan $x<0$ dan ambil $\varepsilon=-x>0$. Ada $N$ dengan $|x_n-x|<\varepsilon$
untuk $n>N$, sehingga khususnya $x_n < x+\varepsilon = 0$. Itu bertentangan dengan
$x_n \ge 0$. Untuk pernyataan kedua, terapkan bagian pertama pada $y_n-x_n \ge 0$,
yang limitnya $y-x$ menurut teorema sebelumnya.
\end{proof}

Perhatikan bahwa kesimpulan yang sah adalah $x \ge 0$, bukan $x>0$: barisan
$1/n$ seluruhnya positif tetapi limitnya nol. Sifat urutan juga tidak boleh dipakai
mundur. Dari $x_n<y_n$ kita hanya memperoleh $x \le y$; kesenjangan ketat bisa
lenyap di limit.

Dengan sifat aljabar, limit barisan rasional dapat dihitung seperti biasa. Untuk
\[
x_n=\frac{3n^2+2n}{n^2-5},
\]
bagi pembilang dan penyebut dengan $n^2$:
$x_n=\dfrac{3+2/n}{1-5/n^2}$. Karena $1/n \to 0$ dan $1/n^2 \to 0$, pembilangnya
menuju $3$ dan penyebutnya menuju $1 \neq 0$, jadi $x_n \to 3$. Nilai numeriknya
memang merayap ke sana: $x_{10}=3.368421$, $x_{100}=3.021511$. Contoh lain,
\[
y_n=\frac{2n^2-n}{3n^2+7} \longrightarrow \frac{2}{3},
\]
dengan $y_{10}=0.618893$ dan $y_{100}=0.663179$; keduanya masih jauh dari $2/3$
pada $n=100$ karena galatnya sekitar $1/n$.

Sifat aljabar hanya berbicara tentang limit yang sudah ada; menerapkannya pada
barisan yang belum diketahui konvergen adalah kekeliruan yang mahal.

\section{Teorema apit dan barisan pembanding}
\tingkat{menengah}

Kadang limit sebuah barisan tidak dapat dihitung langsung, tetapi dapat dikurung
dari atas dan dari bawah oleh dua barisan yang limitnya sama. Kalimat sehari-harinya:
bila dua penjaga berjalan berdekatan sambil mengapit seorang anak, anak itu harus
berhenti di tempat yang sama dengan mereka.

\begin{teorema}[Teorema apit]
\index{teorema apit}
Misalkan $a_n \le x_n \le b_n$ untuk semua $n$ yang cukup besar. Jika
$a_n \to L$ dan $b_n \to L$, maka $x_n \to L$.
\end{teorema}

\begin{proof}[Bukti]
Ambil $\varepsilon>0$. Pilih $N_1$ sehingga $|a_n-L|<\varepsilon$ untuk $n>N_1$,
yakni $a_n>L-\varepsilon$, dan $N_2$ sehingga $b_n<L+\varepsilon$ untuk $n>N_2$.
Ambil juga $N_3$ sehingga pengapitan berlaku untuk $n>N_3$. Untuk
$n>\max(N_1,N_2,N_3)$,
\[
L-\varepsilon < a_n \le x_n \le b_n < L+\varepsilon,
\]
jadi $|x_n-L|<\varepsilon$.
\end{proof}

Sebelum memakainya, kita perlu satu barisan pembanding yang akan muncul terus:
pangkat yang mengecil menuju nol.

\begin{lemma}
Jika $0<q<1$, maka $q^n \to 0$.
\end{lemma}

\begin{proof}[Bukti]
Tulis $1/q=1+h$ dengan $h>0$. Ketaksamaan Bernoulli, yang dibuktikan dengan induksi
di Bab 1, memberi $(1+h)^n \ge 1+nh$. Karena kedua ruas positif,
\[
0 < q^n = \frac{1}{(1+h)^n} \le \frac{1}{1+nh} .
\]
Diberikan $\varepsilon>0$, sifat Archimedes memberi $n$ dengan $1+nh>1/\varepsilon$,
yakni $1/(1+nh)<\varepsilon$. Menurut teorema apit, $q^n \to 0$.
\end{proof}

Sekarang contoh yang tidak dapat diselesaikan dengan membagi pangkat tertinggi.
Untuk $x_n=(2^n+3^n)^{1/n}$, suku terbesar di dalam kurung adalah $3^n$; keluarkan
pangkat itu:
\[
x_n = 3\left(1+\left(\tfrac{2}{3}\right)^{n}\right)^{1/n}.
\]
Karena $(2/3)^n \to 0$ dan $1 \le 1+(2/3)^n \le 2$, kita peroleh
$1 \le x_n/3 \le 2^{1/n}$. Tinggal ditunjukkan $2^{1/n} \to 1$: tulis
$2^{1/n}=1+a_n$ dengan $a_n \ge 0$, maka
$2=(1+a_n)^n \ge 1+na_n$, sehingga $0 \le a_n \le 1/n$. Teorema apit menutup
kedua ujung, dan $x_n \to 3$. Angka Python menegaskan arahnya:
$x_1=5$, $x_2=3.605551$, $x_5=3.075152$, $x_{10}=3.005162$, dan $x_{100}=3$.

\begin{untukanak}
Dua penjaga sekolah berjalan berdampingan menyusuri koridor, dan seorang murid
berjalan di antara mereka. Penjaga di kiri melangkah ke arah pintu keluar, penjaga
di kanan juga, dan keduanya makin rapat sampai bahu mereka hampir bersentuhan.
Murid yang terjepit di tengah tidak perlu tahu jalan; ia hanya perlu tetap di
antara kedua penjaga. Ke mana pun penjaga itu berhenti, di situ pula murid
berhenti.

Sekarang pertanyaan yang membuat penasaran. Bagaimana kalau penjaga kiri berhenti
lebih awal daripada penjaga kanan: masih bisakah kita menebak di mana murid
berhenti? Lalu bagaimana kalau jarak mereka tidak pernah nol, hanya makin kecil
dari hari ke hari? Apakah murid itu tetap punya tempat berhenti yang pasti, atau
terus berjalan tanpa ujung?

Ternyata jawabannya tidak bergantung pada wajah penjaganya, hanya pada dua hal: ke
mana penjaga kiri menuju, dan ke mana penjaga kanan menuju. Kalau keduanya menuju
tempat yang sama, murid di tengah tidak punya pilihan lain. Cara berpikir "terjepit
di antara dua pihak" itulah yang dipakai matematikawan untuk barisan yang rumusnya
terlalu rumit dihitung langsung.
\end{untukanak}

\section{Barisan monoton dan teorema kekonvergenan monoton}
\tingkat{menengah}

Kalimat sehari-harinya: barisan yang hanya bergerak satu arah, dan tidak pernah
melewati pagar tertentu, pasti berhenti di suatu tempat. Inilah tempat kelengkapan
bekerja paling murah, karena limitnya tidak perlu ditebak lebih dahulu.

\kamus{barisan monoton}{monotone sequence}{tangga yang anak tangganya hanya naik,
tidak pernah turun}

Bahasa sehari-harinya sederhana: satu arah saja, tidak pernah berbalik.

\begin{definisi}[Barisan monoton]
Barisan $(x_n)$ disebut \emph{naik}\index{barisan monoton} jika $x_n \le x_{n+1}$
untuk semua $n$, dan \emph{turun} jika $x_n \ge x_{n+1}$ untuk semua $n$.
Barisan disebut monoton jika ia naik atau turun.
\end{definisi}

Gambaran kasarnya: yang melangkah satu arah tanpa melewati pagar pasti berhenti, dan
berhentinya tepat di pagar terdekat yang tidak pernah dilewati.

\begin{teorema}[Kekonvergenan monoton]
\index{teorema kekonvergenan monoton}
Barisan naik yang terbatas di atas konvergen, dan limitnya adalah
$\sup\{x_n : n \in \mathbb{N}\}$.
\end{teorema}

\begin{proof}[Bukti]
Misalkan $(x_n)$ naik dan terbatas di atas. Himpunan sukunya tidak kosong dan
terbatas di atas, jadi menurut sifat supremum di Bab 3 ia mempunyai supremum;
sebut $s=\sup\{x_n\}$. Ambil $\varepsilon>0$. Karena $s-\varepsilon$ bukan batas
atas, ada $N$ dengan $x_N>s-\varepsilon$. Untuk $n>N$, kemonotonan memberi
$x_n \ge x_N>s-\varepsilon$, sedangkan $x_n \le s$ menurut definisi supremum. Jadi
$|x_n-s|<\varepsilon$.
\end{proof}

Akibat langsungnya: untuk barisan seperti itu limitnya dapat dibaca dari himpunan
sukunya, tanpa menghitung suku demi suku.

\begin{akibat}
Barisan turun yang terbatas di bawah konvergen, dan limitnya adalah infimum
himpunan sukunya.
\end{akibat}

\begin{proof}[Bukti]
Terapkan teorema pada $(-x_n)$, yang naik dan terbatas di atas, lalu gunakan sifat
kelipatan skalar pada limit dan kesamaan
$\sup\{-x_n\}=-\inf\{x_n\}$.
\end{proof}

Perhatikan bahwa teorema ini tidak mengatakan limitnya sama dengan batas atas yang
kita pakai. Batas $3$ pada contoh berikut hanya alat bukti; limitnya justru $e$.

Aplikasi pertama adalah barisan rekursif yang membuka bab ini. Misalkan $x_1=1$ dan
$x_{n+1}=\sqrt{2+x_n}$.

Pertama, semua suku berada di antara $1$ dan $2$. Untuk $n=1$ jelas. Bila
$1 \le x_n \le 2$, maka $3 \le 2+x_n \le 4$, sehingga
$1 \le \sqrt{3} \le x_{n+1} \le 2$. Dengan induksi, $1 \le x_n \le 2$ untuk semua
$n$.

Kedua, barisan itu naik. Karena $x_n \le 2$, berlaku
\[
x_{n+1} \ge x_n \iff 2+x_n \ge x_n^2 \iff (2-x_n)(1+x_n) \ge 0,
\]
dan ruas terakhir benar sebab $x_n \le 2$ dan $x_n \ge 1$.

Menurut teorema kekonvergenan monoton, $(x_n)$ konvergen; sebut limitnya $L$.
Karena $x_n \le 2$ untuk semua $n$, sifat urutan memberi $L \le 2$. Sekarang
kuadratkan relasinya: $x_{n+1}^2=2+x_n$. Ruas kiri menuju $L^2$ menurut sifat
perkalian, ruas kanan menuju $2+L$, jadi $L^2=2+L$, atau
$(L-2)(L+1)=0$. Karena $L \ge 1$, satu-satunya kemungkinan adalah $L=2$.

Perhatikan bahwa kita tidak pernah memakai kesinambungan akar; cukup sifat aljabar
limit dan sifat urutan, sehingga peralatan Bab 11 belum diperlukan di sini.

Aplikasi kedua adalah barisan $y_n=(1+1/n)^n$. Ekspansi binomial memberi
\[
y_n = \sum_{k=0}^{n} \binom{n}{k}\frac{1}{n^k}
= \sum_{k=0}^{n} \frac{1}{k!}\prod_{j=1}^{k-1}\left(1-\frac{j}{n}\right),
\]
dengan hasil kali kosong bernilai $1$ untuk $k=0$ dan $k=1$. Setiap faktor
$1-j/n$ naik ketika $n$ bertambah, sehingga setiap suku dengan indeks $k$ naik;
jumlahnya juga mendapat satu suku baru yang positif. Jadi $(y_n)$ naik. Setiap
faktor paling besar $1$, sehingga
\[
y_n \le \sum_{k=0}^{n}\frac{1}{k!} \le 1+1+\frac12+\frac14+\cdots = 3,
\]
karena $1/k! \le 1/2^{k-1}$ untuk $k \ge 1$, yang diperiksa dengan induksi.
Barisan naik dan terbatas di atas oleh $3$, jadi ia konvergen. Limitnya kita beri
nama $e$.

Gambar~\ref{fig:monoton-e} menampilkan barisan itu untuk $n=1$ sampai $n=30$.
Titik-titiknya naik tanpa pernah berbalik, garis putus-putus di $3$ adalah pagar
yang dipakai dalam bukti, dan garis titik-titik di $2.71828$ adalah tempat ia
benar-benar berhenti. Jarak antara pagar dan tujuan tetap ada sepanjang gambar;
pagar bekerja karena tidak pernah dilewati, bukan karena dicapai.

\begin{figure}[htbp]
\centering
\begin{tikzpicture}[x=0.32cm,y=3.0cm,>=Stealth,font=\footnotesize]
  \draw[->,ink] (-0.6,1.85) -- (31.2,1.85) node[right] {$n$};
  \draw[->,ink] (0,1.85) -- (0,3.12) node[above] {$y_n$};
  \draw[dashed,layergrey,thick] (0,3.0) -- (30.6,3.0);
  \node[layergrey,anchor=south] at (20,3.02) {batas bukti $B=3$};
  \draw[densely dotted,layerindigo,very thick] (0,2.71828) -- (30.4,2.71828);
  \node[layerindigo,anchor=west] at (31.3,2.71828) {limit $e\approx2.71828$};
  \foreach \n/\y in {1/2.0000,2/2.2500,3/2.3704,4/2.4414,5/2.4883,6/2.5216,%
      7/2.5465,8/2.5658,9/2.5812,10/2.5937,11/2.6042,12/2.6130,13/2.6206,%
      14/2.6272,15/2.6329,16/2.6379,17/2.6424,18/2.6464,19/2.6500,20/2.6533,%
      21/2.6563,22/2.6590,23/2.6615,24/2.6637,25/2.6658,26/2.6678,27/2.6696,%
      28/2.6713,29/2.6728,30/2.6743}
    \fill[accent] (\n,\y) circle (1.7pt);
  \draw[->,layerteal,thick] (9.5,2.42) -- (15.5,2.60);
  \node[layerteal,anchor=north] at (12,2.40) {naik, tidak pernah berbalik};
\end{tikzpicture}
\caption{Barisan $y_n=(1+1/n)^n$ untuk $n=1$ sampai $n=30$. Titik-titiknya naik dan
tidak pernah berbalik; pagar $B=3$ tidak pernah disentuh; nilai yang dituju adalah
$e\approx2.71828$, bukan pagarnya.}
\label{fig:monoton-e}
\end{figure}

\section{Subbarisan, Bolzano--Weierstrass, dan kriteria Cauchy}
\tingkat{lanjut}

Kalimat sehari-harinya: dari barisan yang terbatas selalu dapat dipetik satu
deretan anak yang berhenti di suatu tempat, walaupun deretan induknya kacau.

\kamus{subbarisan}{subsequence}{menyortir sebagian surat dari tumpukan tanpa
mengubah urutannya}

Subbarisan dibentuk dengan memilih sebagian suku saja: nomor yang dipakai makin
besar, tetapi suku yang dilewati boleh sebanyak apa pun.

\begin{definisi}[Subbarisan]
Misalkan $n_1<n_2<n_3<\cdots$ bilangan asli naik. Barisan $(x_{n_k})_{k \in
\mathbb{N}}$ disebut \emph{subbarisan}\index{subbarisan} dari $(x_n)$.
\end{definisi}

Karena $n_k \ge k$ untuk setiap $k$, yang diperiksa dengan induksi, subbarisan
tidak dapat "lebih lambat" daripada barisan asalnya dalam hal indeks.

Kalimat sehari-harinya: memilih sebagian suku tidak mengubah tempat berhenti, hanya
memperlambat atau mempercepat kedatangannya.

\begin{proposisi}
Jika $x_n \to L$, maka setiap subbarisan $(x_{n_k})$ juga konvergen ke $L$.
\end{proposisi}

\begin{proof}[Bukti]
Ambil $\varepsilon>0$ dan $N$ dengan $|x_n-L|<\varepsilon$ untuk $n>N$. Untuk
$k>N$ berlaku $n_k \ge k>N$, jadi $|x_{n_k}-L|<\varepsilon$.
\end{proof}

Akibatnya sederhana tetapi paling sering dipakai untuk menyangkal: dua tempat
berhenti yang berbeda sudah cukup.

\begin{akibat}
Jika $(x_n)$ mempunyai dua subbarisan dengan limit berbeda, maka $(x_n)$ divergen.
\end{akibat}

\begin{proof}[Bukti]
Bila $(x_n)$ konvergen, semua subbarisannya konvergen ke limit yang sama, dan
ketunggalan limit melarang dua nilai berbeda.
\end{proof}

Akibat ini menjelaskan contoh $(-1)^n$ tanpa argumen kemustahilan: subbarisan
genap menuju $1$, subbarisan ganjil menuju $-1$. Cara kerjanya juga menjelaskan
mengapa divergensi lebih sering dapat dibuktikan daripada kekonvergenan.

Untuk membuktikan Bolzano--Weierstrass, kita butuh lemma berikut. Isinya: setiap
barisan, betapa pun kacau, memuat subbarisan yang monoton.

\begin{lemma}
Setiap barisan real mempunyai subbarisan monoton.
\end{lemma}

\begin{proof}[Bukti]
Sebut indeks $n$ sebuah \emph{puncak} bila $x_n \ge x_m$ untuk semua $m>n$, yakni
bila tidak ada suku sesudahnya yang lebih besar. Ada dua kemungkinan.

Bila puncak tak berhingga banyak, susun $n_1<n_2<\cdots$ sebagai daftar puncak.
Karena $n_1$ puncak, $x_{n_1} \ge x_{n_2}$; karena $n_2$ puncak,
$x_{n_2} \ge x_{n_3}$; dan seterusnya. Hasilnya subbarisan turun.

Bila puncak hanya berhingga, ambil $n_1$ melewati puncak terakhir. Karena $n_1$
bukan puncak, ada $n_2>n_1$ dengan $x_{n_2}>x_{n_1}$. Karena $n_2$ juga melewati
puncak terakhir, ia bukan puncak, jadi ada $n_3>n_2$ dengan $x_{n_3}>x_{n_2}$.
Proses ini berlanjut tanpa henti dan menghasilkan subbarisan naik.
\end{proof}

Isi teorema ini: dari barisan yang terkurung selalu dapat dipetik sebagian suku yang
benar-benar menuju satu titik.

\begin{teorema}[Bolzano--Weierstrass]
\index{bolzano-weierstrass}
Setiap barisan real yang terbatas mempunyai subbarisan yang konvergen.
\end{teorema}

\begin{proof}[Bukti]
Menurut lemma, $(x_n)$ mempunyai subbarisan monoton. Subbarisan dari barisan
terbatas juga terbatas. Barisan monoton yang terbatas konvergen menurut teorema
kekonvergenan monoton. Jadi subbarisan itu konvergen.
\end{proof}

Bukti alternatif memakai pembagian dua interval dan sifat interval bersarang; kedua
bukti bersandar pada kelengkapan, yang tadi dipinjam lewat supremum.

\kamus{barisan Cauchy}{Cauchy sequence}{penumpang yang makin berdesakan di lorong
kereta, tanpa tahu pintu keluarnya di mana}

Sekarang perhatikan masalah yang berbeda. Definisi limit menuntut kita menyebut
$L$ lebih dahulu, padahal dalam banyak perhitungan justru $L$ yang dicari. Yang
kita punya hanya daftar angka. Karena itu kita perlu ukuran kedekatan antar suku
sendiri, tanpa menyebut tujuan.

\begin{definisi}[Barisan Cauchy]
Barisan $(x_n)$ disebut \emph{Cauchy}\index{barisan cauchy} jika untuk setiap
$\varepsilon>0$ terdapat $N$ sehingga untuk semua $m,n>N$ berlaku
$|x_n-x_m|<\varepsilon$.
\end{definisi}

Perhatikan perbedaan urutan gagasannya. Pada definisi limit semua suku dibandingkan
dengan satu bilangan tetap $L$; pada definisi Cauchy suku dibandingkan dengan suku
lain, tanpa bilangan luar. Karena itu definisi Cauchy masuk akal di ruang mana pun
yang punya jarak, sedangkan definisi limit menuntut calon tujuan.

\begin{teorema}[Kriteria Cauchy]
\index{barisan cauchy}
Barisan real konvergen jika dan hanya jika ia Cauchy.
\end{teorema}

\begin{proof}[Bukti]
Misalkan $x_n \to L$ dan ambil $\varepsilon>0$. Ada $N$ dengan
$|x_n-L|<\varepsilon/2$ untuk $n>N$. Untuk $m,n>N$,
\[
|x_n-x_m| \le |x_n-L| + |L-x_m| < \varepsilon .
\]
Jadi konvergen memberi Cauchy.

Sebaliknya, misalkan $(x_n)$ Cauchy. Ambil $\varepsilon=1$; ada $N_0$ sehingga
$|x_n-x_m|<1$ untuk $m,n>N_0$. Dengan memilih $m=N_0+1$ diperoleh
$|x_n|<|x_{N_0+1}|+1$ untuk $n>N_0$, jadi $(x_n)$ terbatas. Menurut
Bolzano--Weierstrass ada subbarisan $x_{n_k} \to L$. Ambil $\varepsilon>0$; pilih
$N_1$ dengan $|x_n-x_m|<\varepsilon/2$ untuk $m,n>N_1$, dan $K$ dengan
$|x_{n_k}-L|<\varepsilon/2$ untuk $k>K$. Pilih $k>K$ yang juga memenuhi
$n_k>N_1$; untuk $n>N_1$,
\[
|x_n-L| \le |x_n-x_{n_k}| + |x_{n_k}-L| < \varepsilon .
\]
Jadi $x_n \to L$.
\end{proof}

Arah kedua memakai Bolzano--Weierstrass, yang memakai teorema kekonvergenan
monoton, yang memakai supremum. Rantai itu tidak dapat diputus: di dalam
$\mathbb{Q}$ teorema ini gagal. Barisan pemotongan desimal
$1;\ 1.4;\ 1.41;\ 1.414;\ 1.4142;\ \dots$ menuju akar dua adalah Cauchy, sebab
selisih dua sukunya paling besar $10^{-k}$ bila keduanya melewati suku ke-$k$,
tetapi tidak ada bilangan rasional yang menjadi limitnya. Bila barisan itu
konvergen ke $L$, maka $a_k^2 \to L^2$; di pihak lain $a_k^2 \to 2$, sehingga
$L^2=2$. Bab 3 sudah menunjukkan tidak ada rasional yang kuadratnya $2$. Jadi
"Cauchy" dan "konvergen" adalah dua sifat berbeda; yang menyamakan keduanya adalah
kelengkapan, dan itulah isi aksioma supremum.

\begin{tantangan}
Barisan $h_n=1+\frac12+\cdots+\frac1n$ tumbuh sangat lambat: $h_{10}=2.928968$,
$h_{100}=5.187378$, $h_{1000}=7.485471$, dan $h_{10^6}=14.392727$. Selisih dua suku
berurutan, yaitu $1/(n+1)$, menuju nol, sehingga barisan ini lolos pemeriksaan
dangkal yang paling sering menipu: suku-sukunya makin rapat. Pertanyaannya, apakah
$(h_n)$ konvergen? Buktikan jawabanmu, dan jelaskan mengapa $h_{10^6}<15$ bukan
bukti bahwa barisan itu terbatas.

Petunjuk 1: jangan periksa suku satu per satu. Jumlahkan dalam blok yang panjangnya
berlipat dua, yaitu satu suku, dua suku, empat suku, delapan suku, dan seterusnya.

Petunjuk 2: blok ke-$k$ berisi $2^{k-1}$ suku, dan suku terkecil di dalamnya paling
kecil $1/2^k$; jadi setiap blok menyumbang sedikitnya $1/2$. Bandingkan hasilnya
dengan $h_{2^m}$ dan lihat apa yang terjadi bila $m$ membesar.

Kalau mentok, kembali ke Teorema Kekonvergenan Monoton: di sana barisan naik
dibuktikan konvergen justru karena ada pagar, dan di sini justru pagarnya yang
tidak ada.
\end{tantangan}

\section{Limit superior dan limit inferior}
\tingkat{lanjut}

Kalimat sehari-harinya: untuk barisan yang berayun, kita masih dapat menyebut
tempat berayun yang paling tinggi dan yang paling rendah, walaupun barisan itu
sendiri tidak punya limit.

\kamus{limit superior}{upper limit}{puncak tertinggi yang masih dikunjungi
barisan sesering apa pun setelah melewati sebarang titik}

Sepanjang bagian ini $(x_n)$ terbatas. Untuk setiap $n$ bentuklah himpunan ekor
$\{x_k : k \ge n\}$ dan tulis
\[
s_n=\sup\{x_k : k \ge n\}, \qquad t_n=\inf\{x_k : k \ge n\}.
\]
Karena ekornya menyusut, $(s_n)$ turun dan $(t_n)$ naik; keduanya terbatas, jadi
keduanya konvergen menurut akibat teorema kekonvergenan monoton.

Kalimat sehari-harinya: yang kita lacak bukan di mana barisan berada, melainkan
seberapa tinggi dan seberapa rendah ia masih berkunjung setelah titik mana pun.

\begin{definisi}[Limit superior dan inferior]
Untuk barisan terbatas $(x_n)$,
\[
\limsup_{n\to\infty} x_n = \lim_{n\to\infty} s_n
\qquad \text{dan} \qquad
\liminf_{n\to\infty} x_n = \lim_{n\to\infty} t_n .
\]
Bilangan pertama disebut \emph{limit superior}\index{limit superior}, yang kedua
\emph{limit inferior}\index{limit inferior}.
\end{definisi}

Karena $s_n \ge t_n$ untuk setiap $n$ dan limit mempertahankan urutan, selalu
berlaku $\liminf x_n \le \limsup x_n$. Keduanya dapat berbeda jauh, dan justru
selisih itulah ukuran kekacauan barisan.

\begin{contoh}
Untuk $x_n=(-1)^n+1/n$, ekornya memuat suku genap $1+1/k$ dan suku ganjil
$-1+1/k$. Supremum ekor ke-$n$ datang dari suku genap terkecil yang masih
tersedia, jadi $s_1=1.5$, $s_2=1.5$, $s_3=1.25$, $s_4=1.25$, $s_5=1.166667$,
dan seterusnya menurun menuju $1$. Infimum ekornya berbeda sifat: suku ganjil
$-1+1/k$ selalu hadir di setiap ekor dan turun menuju $-1$ tanpa pernah
mencapainya, sehingga $t_n=-1$ untuk semua $n$, walaupun tidak ada satu suku pun
yang bernilai $-1$. Maka $\limsup x_n=1$ dan $\liminf x_n=-1$, sedangkan
barisannya divergen.
\end{contoh}

Gambar~\ref{fig:divergen-ayun} memplot suku ke-$n$ terhadap $n$ untuk barisan yang
sama. Deretan titik genap menempel di dekat $1$, deretan titik ganjil menempel di
dekat $-1$, dan jarak kedua kumpulan itu tetap $2$. Tidak ada satu garis mendatar
yang dapat didekati seluruh titik, sejauh apa pun kita berjalan.

\begin{figure}[htbp]
\centering
\begin{tikzpicture}[x=0.62cm,y=1.9cm,>=Stealth,font=\footnotesize]
  \draw[->,ink] (-0.3,0) -- (13.4,0) node[right] {$n$};
  \draw[->,ink] (0,-1.3) -- (0,2.0) node[above] {$x_n$};
  \draw[densely dotted,layerindigo,very thick] (0,1) -- (13.0,1);
  \draw[densely dotted,layerplum,very thick] (0,-1) -- (13.0,-1);
  \node[layerindigo,anchor=south west] at (0.1,1.03) {$\limsup x_n=1$};
  \node[layerplum,anchor=north west] at (0.1,-1.03) {$\liminf x_n=-1$};
  \foreach \n/\y in {1/0.0000,2/1.5000,3/-0.6667,4/1.2500,5/-0.8000,6/1.1667,%
      7/-0.8571,8/1.1250,9/-0.8889,10/1.1000,11/-0.9091,12/1.0833}
    \fill[accent] (\n,\y) circle (2.2pt);
  \draw[<->,layerbrown,thick] (12.5,-1) -- (12.5,1) node[midway,right] {jarak $2$};
\end{tikzpicture}
\caption{Grafik suku ke-$n$ terhadap $n$ untuk $x_n=(-1)^n+1/n$. Titik genap
menempel di dekat $1$, titik ganjil di dekat $-1$, dan jarak kedua kumpulan itu
tetap $2$. Tidak ada pita mendatar yang memuat semua suku, jadi barisan ini
divergen.}
\label{fig:divergen-ayun}
\end{figure}

Kalimat sehari-harinya: nilai tertinggi yang masih dikunjungi barisan sesering apa
pun memang dikunjungi oleh sebagian sukunya, bukan sekadar didekati dari jauh.

\begin{teorema}
Untuk barisan terbatas $(x_n)$, limit superiornya adalah limit subbarisan
terbesar: $\limsup x_n$ sama dengan supremum himpunan semua limit subbarisan yang
konvergen, dan nilai itu sendiri dicapai oleh suatu subbarisan.
\end{teorema}

\begin{proof}[Bukti]
Tulis $S=\limsup x_n$. Pertama, setiap limit subbarisan paling besar $S$. Bila
$x_{n_k} \to L$, maka untuk setiap $k$ berlaku $x_{n_k} \le s_{n_k}$, dan karena
$n_k \ge k$ serta $(s_n)$ turun, $s_{n_k} \le s_k$. Jadi $x_{n_k} \le s_k$, dan
mengambil limit memberi $L \le S$.

Kedua, kita bangun subbarisan yang menuju $S$, mulai dari $n_0=0$. Untuk setiap $k$,
karena $s_n \to S$ dan $(s_n)$ turun, ada $M \ge k$ dengan $s_M<S+1/k$. Ekor
$\{x_n : n \ge M\}$ mempunyai supremum $s_M$ dan memuat indeks sebesar apa pun, jadi
ada $n_k>\max(k,n_{k-1})$ dengan $n_k \ge M$ dan $x_{n_k}>s_M-1/k$. Untuk pilihan itu
\[
S - \frac1k \le s_M - \frac1k < x_{n_k} \le s_M < S + \frac1k ,
\]
sehingga $|x_{n_k}-S|<1/k$. Jadi $x_{n_k} \to S$.
\end{proof}

Akibatnya: kesamaan limit inferior dan limit superior menjadi cara termurah
memeriksa kekonvergenan tanpa menebak limitnya lebih dahulu.

\begin{akibat}
Barisan terbatas $(x_n)$ konvergen ke $L$ jika dan hanya jika
$\liminf x_n = \limsup x_n = L$.
\end{akibat}

\begin{proof}[Bukti]
Bila $x_n \to L$, setiap subbarisan juga menuju $L$, jadi himpunan limit
subbarisan hanya berisi $L$ dan supremumnya $L$; dengan cara yang sama infimumnya
$L$. Sebaliknya, bila kedua nilai sama dengan $L$, maka $t_n \le x_n \le s_n$ dan
kedua pagar menuju $L$, sehingga teorema apit memberi $x_n \to L$.
\end{proof}

Bagi barisan tak terbatas dipakai kesepakatan $\limsup x_n=+\infty$ bila $(x_n)$
tidak terbatas di atas; hukum aritmetika untuk $\infty$ tidak berlaku otomatis, dan
kita tidak memakainya di sini.

\section{Contoh hitung}
\tingkat{menengah}

Kita kembali ke dua barisan pembuka, kali ini dengan angka yang bisa diperiksa
ulang dan dengan langkah antara yang ditulis.

\begin{resep}
\begin{enumerate}
  \item Tulis selisihnya. Bentuk $|x_n-L|$, lalu sederhanakan sampai tinggal satu
        pecahan dalam $n$; untuk contoh di bawah, $|x_n-3|=5/(n+2)$.
  \item Selesaikan untuk $n$. Pecahkan $5/(n+2)<\varepsilon$ menjadi
        $n>5/\varepsilon-2$, dan namai ambang itu $g(\varepsilon)$.
  \item Pilih $N$. Ambil $N=\lfloor g(\varepsilon)\rfloor$, bilangan terkecil yang
        sah, lalu periksa galat pada $n=N$ dan $n=N+1$ sebelum menulis apa pun.
  \item Tutup dengan kalimat definisi: untuk setiap $\varepsilon>0$ dipilih $N$
        itu, maka setiap $n>N$ memenuhi $|x_n-3|<\varepsilon$.
\end{enumerate}
Peringatan urutan: $\varepsilon$ disebut lebih dahulu, $N$ menjawab. Menyusun $N$
lalu mencari $\varepsilon$ yang cocok membalik permainan; hasilnya biasanya $N$
yang bergantung pada $n$, bukan pada $\varepsilon$, dan di titik itu buktinya batal.
\end{resep}

Untuk $x_n=(3n+1)/(n+2)$,
\[
|x_n-3| = \left|\frac{3n+1-3n-6}{n+2}\right| = \frac{5}{n+2}.
\]
Pertidaksamaan $5/(n+2)<\varepsilon$ setara dengan $n>5/\varepsilon-2$. Karena
syaratnya ketat, bilangan bulat terbesar yang masih di bawah ambang itu sudah sah:
$N=\lfloor 5/\varepsilon-2 \rfloor$ adalah pilihan terkecil yang memenuhi definisi.
Untuk $\varepsilon=10^{-2}$ ambangnya $498$ dan $N=498$. Uji batasnya: pada $n=498$
galatnya $5/500=0.01$, yang \emph{tidak} lebih kecil dari $\varepsilon$; pada $n=499$
galatnya $5/501=0.00998004$, baru memenuhi. Untuk $\varepsilon=10^{-4}$ diperoleh
$N=49998$, dengan galat $0.0001$ pada $n=49998$ dan $9.9998 \times 10^{-5}$ pada
$n=49999$.

Nilai sukunya bergerak sangat lambat: $n=10$ memberi $2.583333$, $n=1000$ memberi
$2.995010$, dan galat $10^{-4}$ baru tercapai pada suku ke-$49999$. Empat angka di
depan koma sudah benar jauh sebelum galatnya kecil.

\kertas{Untuk $x_n=(3n+1)/(n+2)$, hitung $N$ yang sah untuk $\varepsilon=10^{-3}$
dengan rumus $N=\lfloor 5/\varepsilon-2 \rfloor$, lalu periksa galat pada $n=N$ dan
$n=N+1$. Jelaskan mengapa galat pada $n=N$ boleh tepat sama dengan $\varepsilon$
sementara definisi tetap terpenuhi.}

Untuk $y_n=(1+1/n)^n$, aritmetika 40 digit memberi galat $0.124539$ pada $n=10$,
$0.00135790$ pada $n=10^3$, dan $1.35914 \times 10^{-7}$ pada $n=10^7$. Hasil kali
$n$ dengan galatnya hampir tetap: $1.245394$, $1.357896$, lalu merayap ke
$1.3591409 = e/2$. Jadi galatnya turun seperti $1/n$, bukan $1/n^2$; satu angka
tambahan di belakang koma menuntut $n$ sepuluh kali lebih besar. Di sinilah tabel
mesin pada pembuka bab rusak: pada $n=10^7$ float64 melaporkan
$1.34327 \times 10^{-7}$, karena galat relatif kecil pada $1+1/n$
terlipatgandakan oleh pangkat $n$.

Untuk rekursi $x_{n+1}=\sqrt{2+x_n}$ dengan $x_1=1$, galatnya terhadap $2$ adalah
$2.679 \times 10^{-1}$ pada $n=2$, $1.711 \times 10^{-2}$ pada $n=4$, dan
$2.615 \times 10^{-7}$ pada $n=12$. Nisbah dua galat berurutan menetap di
$0.250000$.

Angka $1/4$ itu bukan kebetulan. Dari relasinya,
\[
x_{n+1}-2 = \frac{(2+x_n)-4}{\sqrt{2+x_n}+2}
= \frac{x_n-2}{\sqrt{2+x_n}+2},
\]
sehingga
\[
|x_{n+1}-2| = \frac{|x_n-2|}{\sqrt{2+x_n}+2}
\le \frac{|x_n-2|}{2+\sqrt{2}} = 0.292893\,|x_n-2| ,
\]
karena $x_n \ge 1$ memberi $\sqrt{2+x_n} \ge \sqrt{3}$. Batas $0.292893$ berlaku
untuk semua $n$, dan nisbah sebenarnya menuju $1/4$ karena penyebutnya menuju $4$.
Akibat praktisnya: galat $10^{-6}$ tercapai pada $n=12$, $10^{-9}$ pada $n=17$, dan
ketelitian mesin sekitar $n=27$; jauh lebih murah daripada membagi interval dua kali
lipat.

\section{Lab komputasi: empat pemeriksaan}
\tingkat{lanjut}

\kamus{galat pemotongan}{truncation error}{selisih antara potret dan wajah aslinya}

Empat klaim di bab ini dapat diperiksa dengan satu skrip pendek: ambang $N$ yang
eksplisit, laju galat $(1+1/n)^n$, ketelitian mesin, dan nisbah galat rekursi.

\begin{lab}
\begin{verbatim}
import math, mpmath as mp

def N_for(eps):
    N = int(5/eps - 2)
    while 5/(N+2) >= eps: N += 1
    while N > 0 and 5/(N+2) < eps: N -= 1
    return N

print("== 1. N eksplisit untuk (3n+1)/(n+2) -> 3 ==")
for eps in (1e-2, 1e-4):
    N = N_for(eps)
    print("eps=%.0e  N=%d  err(N)=%.6g  err(N+1)=%.6g"
          % (eps, N, abs((3*N+1)/(N+2)-3), abs((3*N+4)/(N+3)-3)))

print("== 2. (1+1/n)^n dalam float64 ==")
for k in range(1, 8):
    n = 10**k
    print("n=1e%d  %.16f  e-y=%.3e" % (k, (1+1/n)**n, math.e-(1+1/n)**n))

print("== 3. galat sebenarnya (mpmath, 40 digit) ==")
mp.mp.dps = 40
for k in (5, 6, 7):
    n = 10**k
    y = (1+mp.mpf(1)/n)**n
    print("n=1e%d  e-y=%s  (float64: %.3e)"
          % (k, mp.nstr(mp.e-y, 6), math.e-(1+1/n)**n))

print("== 4. rekursi x_{n+1}=sqrt(2+x_n), x_1=1 ==")
x, xs = 1.0, []
for _ in range(30):
    xs.append(x); x = math.sqrt(2+x)
for n in (1, 2, 3, 4, 5, 6, 11, 12, 13, 14):
    e = abs(xs[n-1]-2)
    r = "" if n == 1 else "  rasio=%.6f" % (e/abs(xs[n-2]-2))
    print("n=%2d  x=%.12f  galat=%.3e%s" % (n, xs[n-1], e, r))
e6 = next(n for n in range(1, 31) if abs(xs[n-1]-2) < 1e-6)
e9 = next(n for n in range(1, 31) if abs(xs[n-1]-2) < 1e-9)
print("galat < 1e-6 pertama kali pada n=%d ; galat < 1e-9 pada n=%d" % (e6, e9))
\end{verbatim}
\end{lab}

\newpage
Keluarannya, apa adanya:

\begin{verbatim}
== 1. N eksplisit untuk (3n+1)/(n+2) -> 3 ==
eps=1e-02  N=498  err(N)=0.01  err(N+1)=0.00998004
eps=1e-04  N=49998  err(N)=0.0001  err(N+1)=9.9998e-05
== 2. (1+1/n)^n dalam float64 ==
n=1e1  2.5937424601000023  e-y=1.245e-01
n=1e2  2.7048138294215285  e-y=1.347e-02
n=1e3  2.7169239322355936  e-y=1.358e-03
n=1e4  2.7181459268249255  e-y=1.359e-04
n=1e5  2.7182682371922975  e-y=1.359e-05
n=1e6  2.7182804690957534  e-y=1.359e-06
n=1e7  2.7182816941320818  e-y=1.343e-07
== 3. galat sebenarnya (mpmath, 40 digit) ==
n=1e5  e-y=1.35913e-5  (float64: 1.359e-05)
n=1e6  e-y=1.35914e-6  (float64: 1.359e-06)
n=1e7  e-y=1.35914e-7  (float64: 1.343e-07)
== 4. rekursi x_{n+1}=sqrt(2+x_n), x_1=1 ==
n= 1  x=1.000000000000  galat=1.000e+00
n= 2  x=1.732050807569  galat=2.679e-01  rasio=0.267949
n= 3  x=1.931851652578  galat=6.815e-02  rasio=0.254333
n= 4  x=1.982889722748  galat=1.711e-02  rasio=0.251074
n= 5  x=1.995717846477  galat=4.282e-03  rasio=0.250268
n= 6  x=1.998929174953  galat=1.071e-03  rasio=0.250067
n=11  x=1.999998954179  galat=1.046e-06  rasio=0.250000
n=12  x=1.999999738545  galat=2.615e-07  rasio=0.250000
n=13  x=1.999999934636  galat=6.536e-08  rasio=0.250000
n=14  x=1.999999983659  galat=1.634e-08  rasio=0.250000
galat < 1e-6 pertama kali pada n=12 ; galat < 1e-9 pada n=17
\end{verbatim}

Tiga hal yang perlu dibaca dari keluaran itu. Pertama, galat pada $n=N$ tepat sama
dengan $\varepsilon$ di kedua baris; itulah sebabnya definisi memakai syarat ketat
$n>N$. Kedua, kolom $e-y$ berlipat sepersepuluh setiap kali $n$ berlipat sepuluh
sampai $n=10^6$, lalu berhenti mengikuti pola; bagian 3 menunjukkan bahwa yang
berhenti adalah ketelitian float64, bukan kekonvergenan barisan. Galat pemotongan
terus mengecil seperti $e/(2n)$, tetapi galat pembulatan mengambil alih begitu
keduanya bertemu. Ketiga, nisbah galat pada bagian 4 turun ke $0.250000$ dan tinggal
di sana, tepat seperti perhitungan tangan sebelumnya. Program tidak membuktikan apa
pun: ia hanya memeriksa klaim yang sudah kita buktikan.

\begin{untukdosen}
Urutan mengajar: mulai dari permainan $\varepsilon$-$N$ di kelas, bukan dari
teorema. Tunjuk dua baris mahasiswa; satu menyebut $\varepsilon$, satu menjawab $N$
untuk $x_n=(3n+1)/(n+2)$, dan dosen mencatat sampai rumusnya muncul sendiri.
Sesudah itu ketunggalan, sifat aljabar, teorema apit, barisan monoton, lalu
subbarisan dan Cauchy. Limit superior ditaruh di akhir, sebagai "limit untuk
barisan yang tidak punya limit".

Dua miskonsepsi khas. Pertama, mahasiswa menganggap $N$ boleh bergantung pada $n$,
lalu menulis "untuk setiap $n$ ada $N=n$" dan menyangka definisi terpenuhi;
gejalanya tampak ketika mereka tidak dapat menghasilkan satu $N$ untuk seluruh
ekor. Kedua, mereka menyamakan "suku makin rapat" dengan "konvergen" dan memakai
kriteria Cauchy tanpa memeriksa ruangnya; pemotongan desimal akar dua di dalam
$\mathbb{Q}$ adalah penawar yang efektif.

Aktivitas papan tulis: minta mahasiswa menuliskan negasi definisi limit, lalu
buktikan divergensi $(-1)^n$ dengan negasi itu, bukan dengan kontradiksi. Kelas
siap lanjut bila mereka dapat menghasilkan $N$ eksplisit tanpa diberi rumus dan
menyebut satu barisan terbatas yang divergen.
\end{untukdosen}

\section{Ringkasan}
\tingkat{dasar}

\begin{ringkasan}
\begin{itemize}
  \item $x_n \to L$ berarti untuk setiap $\varepsilon>0$ ada $N$ sehingga
        $|x_n-L|<\varepsilon$ untuk semua $n>N$. $\varepsilon$ dipilih lebih
        dahulu; $N$ menjawab, dan boleh bergantung pada $\varepsilon$ saja.
  \item Limit tunggal, dan barisan konvergen terbatas. Kebalikannya tidak berlaku:
        keterbatasan hanya syarat perlu.
  \item Limit menghormati $+$, $\times$, dan pembagian dengan penyebut tak nol;
        urutan tidak pernah dilompati, hanya dipertahankan lemah.
  \item Teorema apit mengubah pengurungan menjadi limit, dan lemma $q^n \to 0$
        memberi pembanding yang dipakai berulang. Barisan monoton yang terbatas
        selalu konvergen, dan limitnya supremum atau infimum.
  \item Setiap barisan terbatas punya subbarisan konvergen; barisan real konvergen
        jika dan hanya jika Cauchy, dan kesetaraan itu adalah kelengkapan.
  \item $\liminf \le \limsup$, dan keduanya sama dengan $L$ tepat ketika barisan
        konvergen ke $L$.
\end{itemize}
\end{ringkasan}

\section{Latihan}
\tingkat{menengah}

\begin{latihan}
\begin{enumerate}
  \item Buktikan dengan definisi bahwa $x_n=(2n-1)/(3n+4)$ konvergen ke $2/3$, dan
        berikan $N$ eksplisit sebagai fungsi $\varepsilon$. Bandingkan dengan
        ambang $n$ terkecil yang memenuhi.
  \item Buktikan: jika $x_n \to L$ maka $|x_n| \to |L|$. Tunjukkan dengan contoh
        bahwa kebalikannya salah.
  \item Hitung $\limsup$ dan $\liminf$ dari $x_n=1+(-1)^n+1/n$ dan dari
        $y_n=(-1)^n(1+1/n)$. Periksa jawabanmu dengan Python untuk $n$ sampai
        $1000$.
  \item Nilai klaim ini: "Karena $(1+1/n)^n$ naik dan terbatas di atas oleh $3$,
        limitnya adalah $3$." Benar, salah, atau belum cukup bukti? Sebutkan
        bagian mana yang benar dan bagian mana yang keliru.
  \item Buktikan bahwa $(2^n+3^n)^{1/n} \to 3$ dengan menuliskan kedua pagar
        secara lengkap, termasuk bukti $2^{1/n} \to 1$.
  \item Untuk rekursi $x_{n+1}=\sqrt{2+x_n}$ dengan $x_1=1$, hitung dengan Python
        suku terkecil dengan galat di bawah $10^{-6}$ dan $10^{-12}$. Bandingkan
        dengan tebakan dari nisbah $1/4$.
  \item Buktikan bahwa barisan Cauchy terbatas, tanpa memakai Bolzano--Weierstrass.
        Lalu buktikan bahwa barisan monoton yang terbatas adalah Cauchy.
  \item Untuk $x_n=(3n+1)/(n+2)$ dan $\varepsilon=10^{-3}$, tentukan $N$ terkecil
        yang sah, lalu hitung galat pada $n=N$, $N+1$, dan $N+2$ dengan Python.
        Jelaskan mengapa $N-1$ tidak sah.
  \item Buktikan: jika setiap subbarisan dari $(x_n)$ mempunyai subbarisan yang
        konvergen ke $L$, maka $x_n \to L$.
\end{enumerate}

Petunjuk untuk tiga butir tersulit. Butir 5: gunakan
$1 \le (1+(2/3)^n)^{1/n} \le 2^{1/n}$ lalu lemma $2^{1/n} \to 1$ lewat
ketaksamaan Bernoulli. Butir 7: untuk bagian kedua, bandingkan $|x_n-x_m|$ dengan
$x_m-x_n$ ketika $m>n$ dan gunakan kekonvergenan monoton. Butir 9: andaikan
$x_n$ tidak menuju $L$, bentuk subbarisan yang berjarak tetap dari $L$, lalu pakai
Bolzano--Weierstrass.
\end{latihan}

\begin{untukpeneliti}
Yang benar-benar terbukti di sini hanyalah kesetaraan berikut, dan semuanya
bersandar pada satu aksioma: sifat supremum. Dalam lapangan terurut, pernyataan
berikut ekuivalen: setiap himpunan tak kosong yang terbatas di atas punya
supremum; setiap barisan naik yang terbatas konvergen; setiap barisan Cauchy
konvergen; setiap barisan terbatas punya subbarisan konvergen. Rantai buktinya ada
di Rudin, \emph{Principles of Mathematical Analysis}, Bab 3 (bagian Convergent
Sequences, Subsequences, Cauchy Sequences) dan Abbott, \emph{Understanding
Analysis}, Bab 2 (bagian 2.2 sampai 2.6). Yang populer tetapi tidak terbukti:
"barisan Cauchy selalu konvergen". Itu salah tanpa kelengkapan, dan contoh
pemotongan desimal akar dua di dalam $\mathbb{Q}$ adalah penyangkalnya.

Pertanyaan yang masih terbuka: seberapa cepat sebuah barisan konvergen, dan apakah
laju itu selalu dapat dihitung. Jawabannya tidak. Hasil klasik Specker (1949)
memberi barisan rasional yang dapat dihitung, naik, dan terbatas, tetapi limitnya
tidak dapat dihitung; jadi "konvergen" tidak otomatis memberi batas galat yang bisa
diprogram.

Batas keberlakuan: seluruh bab ini bekerja di $\mathbb{R}$ dengan jarak biasa.
Di ruang metrik umum, Bolzano--Weierstrass berubah menjadi pernyataan tentang
kekompakan, dan kriteria Cauchy tetap menjadi syarat perlu tetapi belum cukup.
Teorema apit juga diam bila kedua pagar tidak menuju nilai yang sama; di situ kita
butuh pembanding yang lebih halus, dan itulah bahan Bab 5.
\end{untukpeneliti}

\peta{Bab 4 pada Rudin, \emph{Principles of Mathematical Analysis}: Bab 3 (Numerical
Sequences and Series, dengan bagian Convergent Sequences, Subsequences, Cauchy
Sequences, Upper and Lower Limits, dan Some Special Sequences). Pada Abbott,
\emph{Understanding Analysis}: Bab 2 (The Limit of a Sequence, The Algebraic and
Order Limit Theorems, The Monotone Convergence Theorem and a First Look at Infinite
Series, Subsequences and the Bolzano--Weierstrass Theorem, The Cauchy Criterion).
Aksioma dan akibat kelengkapan yang dipakai di sini dibangun di Bab 3 buku ini,
bukan diulang.}
